\chapter{Marco Teórico y Estado del Arte}

\section{Teorías del Comportamiento Electoral}
En esta sección se analiza la evolución de las teorías del comportamiento político, desde el enfoque sociológico de la Escuela de Columbia hasta los modelos de elección racional. Esta base teórica es fundamental para justificar el uso de indicadores macroeconómicos como predictores de la intención de voto.

\subsection{El Modelo Sociológico y la Identidad de Grupo}
\subsection{El Voto Económico: Teoría de la Recompensa y el Castigo}

\section{Evolución de la Predicción Electoral}
Se examina la transición entre los métodos cualitativos y los modelos cuantitativos basados en datos estructurales.

\subsection{Limitaciones de la Metodología Basada en Encuestas}
\subsection{Modelos Econométricos y de ``Fundamentales''}

\section{Data Science Aplicado a la Ciencia Política}
Estado actual de la integración de algoritmos de aprendizaje supervisado en el análisis de grandes volúmenes de datos sociales.

\subsection{Algoritmos de Gradient Boosting en Series Temporales Sociales}
\subsection{Análisis Espacial y Teoría de Grafos en Sociología}

\section{El Sistema Electoral Español}
Análisis técnico de la Ley D'Hondt y el impacto de la fragmentación del sistema de partidos tras el fin del bipartidismo en 2015.
